\section{Conclusion and future work}\label{sec:conclusion}
We introduced an auditable separation between low-storage anomaly ranking and few-shot alarm reliability. A target-only alarm built from $k$ calibration scores cannot fire below $\alpha=1/(k{+}1)$, and corruption can make its observed false-alarm behavior anti-conservative. Source normals can produce useful sub-floor empirical operating points, as the historical SC3R diagnostics show, but the strict nested experiment establishes a sharper boundary: category-level Hoeffding certification returns the zero threshold in all 960 gate cells across MVTec, VisA, and two external-transfer directions. The accompanying feasibility result explains why---the available three or four certification categories are orders of magnitude below what the declared simultaneous levels require. Image-level fixed-archive certificates are sometimes nonzero but cannot replace that failed estimand. CRR should therefore be read as a reliability audit and protocol framework, not as an AUROC state-of-the-art, adversarial-robustness, or certified target-control claim.

Future work should prioritize category-rich source archives or a pre-specified tighter category-risk construction with demonstrable finite-sample validity; adding seeds or source images cannot solve a shortage of independent categories. A separately evaluated condition detector and an audited vision--language baseline remain secondary extensions.
